\section{関連研究}
\label{sec:related_work}

\paragraph{教師信号としての3D Gaussianシーン。}
3D Gaussian Splatting (3DGS)は、異方性Gaussianによって放射場を表現し、
高品質かつリアルタイムな新規視点レンダリングを可能にする
\citep{kerbl2023gaussian}。Synthetic-to-realの研究では、実画像が生成される分布に
含まれることを期待して、シミュレータ全体をランダム化することが一般的である
\citep{tobin2017domain}。本研究の利用法はこれとは異なる。Gaussian splatは対象の視覚ドメイン
から再構成され、制御する変数はカメラ姿勢、対象コート、およびビュー軌道である。
さらに重要なことに、レンダラは計量的なコートアライメントの下流に位置するため、
画像のレンダリングに用いたものと同じ変換から、その幾何ラベルも生成される。

\paragraph{スポーツフィールドの位置合わせ。}
スポーツフィールドの位置推定には、検出されたフィールドマーキング、構造化マッチング、
ホモグラフィ推定、およびカメラキャリブレーションが用いられてきた
\citep{homayounfar2017sports}。近年のプロトコルでは、キャリブレーション品質をより
慎重に測定し、可観測性を高めるために点とラインを組み合わせている
\citep{magera2024procc,gutierrez2026pnlcalib}。これらの手法は物理的なライン
トポロジーを用いる根拠となるが、本研究の学習対象は制約のないホモグラフィではない。
本研究では、相互に逆となる計量的なシーン--コート変換を公開し、カメラ内部パラメータを
保持し、新たなキャリブレーション済みビューを生成する。分離したホールドアウト分割に
より、固定されたアライメントがデータセットの正規情報となる前に検証する。

\paragraph{基盤特徴とグローバルクエリ。}
自己教師ありvision Transformerは、密予測タスクへ良好に転移するパッチ特徴を提供する。
本研究では、固定またはfine-tune可能な画像バックボーンとして\DINO{}を用いる
\citep{simeoni2025dinov3}。DETRのようなオブジェクトクエリアーキテクチャは、学習可能な
トークンが可変個の空間トークン集合を一つの大域予測へ集約できることを示している
\citep{carion2020detr}。Rotary position encodingは、空間的関係をアテンションへ直接
導入する\citep{su2021roformer}。本研究のコートクエリは意図的に位置符号を持たせず、
パッチのquery/keyベクトルには2D rotary encodingを適用する。出力は交換可能な
オブジェクトの集合ではなく、単一のカメラ--コート仮説である。

\paragraph{微分可能な幾何整合性。}
微分可能なサンプリングと投影は、画像予測と幾何パラメータを長らく結び付けてきた
\citep{jaderberg2015spatial}。Soft-argmaxおよび積分座標による定式化は、勾配を
途切れさせることなくヒートマップを連続座標へ変換する
\citep{luvizon2019human}。連続な6D回転表現は、回転の直接回帰を困難にする不連続性を
回避する\citep{zhou2019continuity}。本研究では、これらの要素を厳密な意味論的契約の
下で組み合わせる。すなわち、各ヒートマップチャネルは厳密に一つの物理的コート点に
対応し、パディング画素には確率質量を与えず、有効な教師マスクは生成された正解データ
のみから定まり、予測深度が負の場合には暗黙に除外せず罰則を科す。
